\section{Data, Simulation and Estimators}
\label{sec:method}

\textbf{Real drives and labels.}
LiRA-CD~\cite{skar2023lira,lira_data} contains Renault Zoe cars instrumented with an AutoPi telematics unit (three-axis accelerometer, GPS, CAN bus) fixed to the frame in the middle console near the front axle~\cite{skar2023lira}, and P79 laser profiles of the same lanes.
We use the four routes with car data, two urban roads in Copenhagen (CPH1, CPH6) and two motorways (M3, M13), in both directions; the accelerometer is logged at \liraRateLow{} or \liraRateHigh\,Hz and speed comes from the CAN bus.
After map matching to the P79 chainage, the data contain \liraUnits{} segment-passes over \liraSegments{} segments of 100\,m (\liraKm\,km), about ten passes per segment.
Labels are the IRI per 100\,m (mean of both wheel paths), recomputed from the P79 profiles with an implementation of the quarter-car algorithm that reproduces the profilometer's own IRI ($r=$\,\iriValR).
Urban segments average \liraIriUrban\,m/km at \liraSpeedUrban\,km/h, motorway segments \liraIriMotorway\,m/km at \liraSpeedMotorway\,km/h; the surveys preceded the drives by two to fifteen months.

\textbf{Simulator.}
A linear four-degree-of-freedom half-car (bounce, pitch, two unsprung masses) is driven by the mean of the two wheel-path profiles under the front axle and the same profile delayed by the wheelbase under the rear axle, after a 250\,mm tyre-envelope average; simulated units are labelled like the real ones.
Vehicle parameters are drawn from passenger-car ranges~\cite{gillespie1992} (sprung mass 1000--1800\,kg, body frequency 1.0--1.6\,Hz, damping ratio 0.20--0.45) and speed independently of roughness (20--110\,km/h).
A sensor model adds an optional holder resonance, anti-aliasing, a mounting rotation, irregular sampling, scale and bias errors, noise, quantisation and sample loss.
The reference simulator \textbf{V0}, fixed before any test evaluation, uses 30 random vehicles over the measured profiles of the training routes and a 50/50 mixture of a telematics-like sensor (55--62\,Hz, 0.01\,g, rigid) and a phone-like sensor (90--210\,Hz) with a holder resonance at 15--40\,Hz.
Eight variants change one ingredient: either half of V0's sensor mixture alone (\textbf{V1} telematics-like, \textbf{V2} phone-like), the missing cells of a sensor-family $\times$ holder design (\textbf{V7}, \textbf{V8}), an ideal sensor (\textbf{V3}), a telematics sensor matched to statistics of the LiRA-CD streams alone (\textbf{V5}) or mixed with the phone-like sensor (\textbf{V6}), and one mid-range vehicle instead of 30 (\textbf{V4}); \textbf{V9} replaces the measured by synthetic ISO~8608 profiles~\cite{iso8608}.
Simulated drives never use the profile of the test route.

\textbf{Features and estimators.}
Real and simulated accelerometer streams pass the same code: resampling to 200\,Hz, zero-phase 20\,Hz low-pass, decimation to 50\,Hz, projection on the pass-mean gravity direction (orientation invariance) and a 0.3\,Hz high-pass.
A unit is one pass over one 100\,m segment, described by its mean speed and speed variation and by 11 features of the dynamic vertical acceleration: log RMS, peak statistics, crest factor, kurtosis, RMS in the 0.5--2, 2--5, 5--10 and 10--15\,Hz bands, above 1\,Hz and at wavelengths below 5\,m.
The estimators predict log IRI: a road-class prior (map only); LightGBM~\cite{ke2017lightgbm} on speed only; a vibration index $\log\mathrm{IRI}=a+b\log(\mathrm{RMS}/v^{\alpha})$; a calibrated quarter-car inversion, which fits a quarter car to the car's empirical transmissibility on the training routes and inverts each unit's spectrum through the Golden Car (it reproduces the IRI of true profile spectra with $r=$\,\invCheckR); LightGBM on all features; an MLP; and a 1D-CNN on the resampled acceleration (five-seed ensembles).
The untrained statistic RMS$/v^{0.5}$ serves as a ranking baseline.
Simulation-trained (zero-shot) estimators are fitted to the simulated units of one variant.
LightGBM trained on whole real training routes selects one of six configurations by inner leave-one-route-out validation; all other LightGBM models use the first configuration.

\textbf{Paired data.}
Paired data are $k$\,km of training-route road, drawn as contiguous 1\,km blocks (20 draws per $k$, shared by all methods), with all real passes over them: at $k=5$, \kdUnitsFive{} segment-passes on average, about \kdDriveKmFive\,km of driving over \kdKmFive\,km of road.
To test whether this density is needed, each segment also keeps at most $m\in\{1,2,5\}$ passes.
The paired passes are combined with LightGBM trained on simulated drives over all training-route profiles by a level offset (median log ratio of label to prediction), an affine map in log space, or joint training on simulated and real units with equal total weight; the comparators are LightGBM and the vibration index trained on the paired passes alone.

\textbf{Evaluation.}
Each route is held out in turn, both directions together, and nothing from it enters training, calibration or model selection.
A segment is predicted by the mean of its pass predictions.
We report MAE, bias and Spearman's $\rho$; \emph{within-route} values average the four routes, \emph{pooled} values cover all \liraSegments{} segments, and the \emph{level} of a road class is the mean log ratio of prediction to label over its segments, averaged over its two routes.
Segment errors are spatially correlated (autocorrelation \acfGbmHundred{} at 100\,m), so paired differences are bootstrapped over 5\,km blocks (40 clusters; single-estimator spectra and elasticities: 1\,km blocks).
The estimators, the simulator V0 and the evaluation were fixed before any test result; the variant, calibration and spectral analyses were designed after the test routes had been evaluated and are exploratory.
